\subsection{\texorpdfstring{\(\mathfrak{S}\)-CONVERGENCE}{S-CONVERGENCE}}

DEFINITION 2. Let X be a set, Y a uniform space, S a set of subsets of X. The uniformity of uniform convergence in the sets of S, or simply the uniformity of S-convergence, is the coarsest uniformity on F(X; Y) which makes uniformly continuous the restriction mappings \(u \to u|A\) of F(X, Y) into the uniform spaces \(\mathcal{F}_u(A; Y)\) , where A runs through S. The uniform space obtained by endowing F(X; Y) with the uniformity of S-convergence is denoted by \(\mathcal{F}_{\mathfrak{S}}(X; Y)\) .

The topology induced by the uniformity of \(\mathfrak{S}\) -convergence is called the topology of \(\mathfrak{S}\) -convergence; it is the coarsest for which all the mappings \(u \to u|A\) of \(\mathcal{F}(X; Y)\) into \(\mathcal{F}_u(A; Y) (A \in \mathfrak{S})\) are continuous (Chapter II, § 2, no. 3, Proposition 4, Corollary).

A filter \(\Phi\) on \(\mathcal{F}(\mathbf{X};\mathbf{Y})\) converges to \(u_{0}\) with respect to the topology of \(\mathfrak{S}\) -convergence if and only if \(u|\mathrm{A}\) converges uniformly to \(u_{0}|\mathrm{A}\) with respect to \(\Phi\) for all \(\mathrm{A}\in \mathfrak{S}\) (Chapter I, § 7, no. 6, Proposition 10), and \(\Phi\) is therefore said to converge uniformly to \(u_{0}\) on the sets of \(\mathfrak{S}\) .

A filter base \(\Phi\) on \(\mathcal{F}_{\mathfrak{S}}(\mathbf{X};\mathbf{Y})\) is a Cauchy filter base if and only if, for each \(\mathbf{A}\in\mathfrak{S}\) , the image of \(\Phi\) under the mapping \(u\to u|\mathbf{A}\) is a Cauchy filter base on \(\mathcal{F}_{u}(\mathbf{A};\mathbf{Y})\) (Chapter II, § 3, no. 1, Proposition 4).

Let \(f\) be a mapping of a topological (resp. uniform) space \(Z\) into \(\mathcal{F}_{\mathfrak{S}}(X; Y)\) . Then \(f\) is continuous (resp. uniformly continuous) if and only if, for each \(A \in \mathfrak{S}\) , the mapping \(z \to f(z)|A\) of \(Z\) into \(\mathcal{F}_u(A; Y)\) is continuous (resp. uniformly continuous) (Chapter I, § 2, no. 3, Proposition 4; Chapter II, § 2, no. 3, Proposition 4).

Finally, let \(\mathbf{M}\) be a subset of \(\mathcal{F}_{\mathfrak{S}}(\mathbf{X};\mathbf{Y})\) ; then \(\mathbf{M}\) is precompact if and only if, for each \(\mathbf{A} \in \mathfrak{S}\) , the set of restrictions \(u|\mathbf{A}\) for \(u \in \mathbf{M}\) is a precompact subset of \(\mathcal{F}_u(\mathbf{A};\mathbf{Y})\) (Chapter II, § 4, no. 2, Proposition 3).

Remarks. 1) The general definition of the entourages of an initial uniformity (Chapter II, § 2, no. 3, Proposition 4) shows that a fundamental system of entourages of \(\mathcal{F}_{\mathfrak{S}}(\mathrm{X};\mathrm{Y})\) may be obtained as follows: for each \(\mathbf{A}\in \mathfrak{S}\) and each entourage \(\mathbf{V}\) of a fundamental system of entourages \(\mathfrak{B}\) of \(\mathbf{Y}\) , let \(W(\mathbf{A},\mathbf{V})\) be the set of all pairs of mappings \((u,v)\) of \(\mathbf{X}\) into \(\mathbf{Y}\) such that \((u(x),v(x))\in \mathbf{V}\) for each \(x\in \mathbf{A}\) ; as \(\mathbf{A}\) runs through \(\mathfrak{S}\) and \(\mathbf{V}\) runs through \(\mathfrak{B}\) , the finite intersections of the sets \(W(\mathbf{A},\mathbf{V})\) form a fundamental system of entourages of \(\mathcal{F}_{\mathfrak{S}}(\mathrm{X};\mathrm{Y})\) .

This description shows immediately that if \(\mathfrak{S},\mathfrak{S}'\) are two sets of subsets of \(\mathbf{X}\) such that \(\mathfrak{S} \subset \mathfrak{S}'\) , then the uniformity of \(\mathfrak{S}'\) -convergence is finer than that of \(\mathfrak{S}\) -convergence.

\begin{enumerate}
\def\labelenumi{\arabic{enumi})}
\setcounter{enumi}{1}
\tightlist
\item
  However, the uniformity of \(\mathfrak{S}\) -convergence is unaltered by replacing \(\mathfrak{S}\) by the set \(\mathfrak{S}'\) of all subsets of X which are contained in finite unions of sets of \(\mathfrak{S}\) . In the study of \(\mathfrak{S}\) -convergence we may therefore always restrict ourselves to the case where the set \(\mathfrak{S}\) satisfies the following two conditions:
\end{enumerate}

\((\mathbf{F}_{\mathbf{I}}^{\prime})\) Every subset of a set of \(\mathfrak{S}\) belongs to \(\mathfrak{S}\) .

\((F_{\mathrm{II}}^{\prime})\) Every finite union of sets of \(\mathfrak{S}\) belongs to \(\mathfrak{S}\) .

If \((\mathbf{F}_{\mathrm{II}}^{\prime})\) is satisfied, we obtain a fundamental system of entourages of \(\mathcal{F}_{\mathfrak{S}}(\mathbf{X};\mathbf{Y})\) by taking all the sets \(\boldsymbol{W}(\mathrm{A},\mathrm{V})\) , where A runs through S and V runs through a fundamental system of entourages of Y.

\begin{enumerate}
\def\labelenumi{\arabic{enumi})}
\setcounter{enumi}{2}
\tightlist
\item
  The uniformity of \(\mathfrak{S}\) -convergence is the inverse image, under the mapping \(u \to (u|\mathbf{A})_{\mathbf{A} \in \mathfrak{S}}\) of \(\mathcal{F}(\mathbf{X}; \mathbf{Y})\) into \(\prod_{\mathbf{A} \in \mathfrak{S}} \mathcal{F}_u(\mathbf{A}; \mathbf{Y})\) , of the uniformity of this product space (Chapter II, § 2, no. 6, Proposition 8). If \(\mathfrak{S}\) is a covering of \(\mathbf{X}\) , this mapping is injective and \(\mathcal{F}_{\mathfrak{S}}(\mathbf{X}; \mathbf{Y})\) is therefore isomorphic to the uniform subspace of \(\prod_{\mathbf{A} \in \mathfrak{S}} \mathcal{F}_u(\mathbf{A}; \mathbf{Y})\) which is the image of this mapping.
\end{enumerate}

PROPOSITION 1. If Y is Hausdorff and S is a covering of X, then the space \(\mathcal{F}_{\mathfrak{S}}(X;Y)\) is Hausdorff.

Let \(u, v\) be two elements of \(\mathcal{F}_{\mathfrak{S}}(\mathrm{X}; \mathrm{Y})\) such that \((u, v) \in W(\mathrm{A}, \mathrm{V})\) for every entourage \(\mathrm{V}\) of \(\mathrm{Y}\) and every \(\mathrm{A} \in \mathfrak{S}\) . Since \(\mathrm{Y}\) is Hausdorff it follows that \(u\) and \(v\) coincide on every set \(\mathrm{A} \in \mathfrak{S}\) , and since \(\mathfrak{S}\) covers \(\mathrm{X}\) we must have \(u = v\) .

Remarks. 4) Let H be a subset of \(\mathcal{F}(\mathbf{X};\mathbf{Y})\) . By abuse of language, the uniformity (resp. topology) induced on H by the uniformity (resp. topology) of \(\mathfrak{S}\) -convergence on \(\mathcal{F}(\mathbf{X};\mathbf{Y})\) is called the uniformity (resp. topology) of \(\mathfrak{S}\) -convergence on the set H.

\begin{enumerate}
\def\labelenumi{\arabic{enumi})}
\setcounter{enumi}{4}
\tightlist
\item
  Let L be a set filtered by a filter \(\mathfrak{G}\) , and let \(\lambda \to u_{\lambda}\) be a mapping of L into \(\mathcal{F}_{\mathfrak{S}}(\mathbf{X};\mathbf{Y})\) which has a limit \(v\) with respect to \(\mathfrak{G}\) ; we say then that, with respect to the filter \(\mathfrak{G}\) , the mappings \(u_{\lambda}\) of X into Y converge uniformly to \(v\) {[}or that the family \((u_{\lambda})\) is uniformly convergent to \(v\) {]} in every set of \(\mathfrak{S}\) . If \(\mathbf{L} = \mathbf{N}\) and \(\mathfrak{G}\) is the Fréchet filter, we omit mention of \(\mathfrak{G}\) in this statement.
\end{enumerate}

More particularly, suppose that there is a commutative and associative law of composition (written additively) defined on Y. If \((u_{n})\) is any sequence of mappings of X into Y, let \(v_{n}\) be the mapping defined by

\[
v _ {n} (x) = \sum_ {k = 0} ^ {n} u _ {k} (x) \quad (n \in \mathbf {N});
\]

we say that the series whose general term is \(u_n\) is uniformly convergent in every set of \(\mathfrak{S}\) if the sequence \((v_n)\) is uniformly convergent in every set of \(\mathfrak{S}\) . Likewise we define a uniformly summable family \((u_\lambda)_{\lambda \in \mathbf{L}}\) of mappings of \(\mathbf{X}\) into \(\mathbf{Y}\) by considering the mappings \(x \to \sum_{\lambda \in \mathbf{J}} u_\lambda(x)\) for all finite subsets

J of L and the limit of these mappings in \(\mathcal{F}_{\mathfrak{S}}(\mathbf{X};\mathbf{Y})\) with respect to the directed set of finite subsets of L (Chapter III, § 5, no. 1).

\begin{enumerate}
\def\labelenumi{\arabic{enumi})}
\setcounter{enumi}{5}
\tightlist
\item
  It follows immediately from Definitions 1 and 2 that, for every \(x \in \bigcup_{\mathbf{A} \in \mathfrak{S}} \mathbf{A}\) , the mapping \(u \to v(x)\) of \(\mathcal{F}_{\mathfrak{S}}(\mathbf{X}; \mathbf{Y})\) into Y is uniformly continuous. Hence, in particular, if \(\overline{\mathbf{H}}\) denotes the closure of a subset H of \(\mathcal{F}_{\mathfrak{S}}(\mathbf{X}; \mathbf{Y})\) , we have \(\overline{\mathbf{H}}(x) \subset \overline{\mathbf{H}(x)}\) for all \(x \in \bigcup_{\mathbf{A} \in \mathfrak{S}} \mathbf{A}\) (Chapter I, § 2, no. 1, Theorem 1).
\end{enumerate}

